\metatitle{Coxeter groups and types}
\metadescription{Coxeter systems, Weyl groups, classical types, and a
translation table between type A, B, C, D, and affine terminology.}
\metakeywords{Coxeter group,Weyl group,type A,type B,type C,type D,affine type,
Bruhat order}

\section[coxeterGroups]{Coxeter groups and types}

A \defin{Coxeter system}\label{coxeterSystemDefinition} is a pair \((W,S)\)
where \(W\) is generated by a set \(S\) of simple reflections and has a
presentation
\[
  W=\langle s\in S : (st)^{m_{st}}=e\rangle,
\]
with \(m_{ss}=1\), \(m_{st}=m_{ts}\in\{2,3,\dotsc\}\cup\{\infty\}\), and no
relation imposed when \(m_{st}=\infty\).  The basic combinatorics of \(W\)
is controlled by \hyperref[reducedWord]{reduced words},
\hyperref[bruhatOrder]{Bruhat order}, weak order, descents, and parabolic
subgroups.  A standard reference is \cite{Bjorner2005}.

Finite Weyl groups are Coxeter groups coming from
\hyperref[rootSystem]{root systems}.  Every Weyl group is a finite Coxeter
group, but not every finite Coxeter group is crystallographic.

\begin{example}
The symmetric group \(\symS_n\) is the Coxeter group of type \(A_{n-1}\).
The simple reflections are the adjacent transpositions
\[
  s_i=(i,i+1),\qquad 1\leq i\leq n-1.
\]
They satisfy \(s_i^2=e\), the braid relation
\[
  s_is_{i+1}s_i=s_{i+1}s_is_{i+1},
\]
and the commutation relation \(s_is_j=s_js_i\) whenever \(|i-j|\geq 2\).
\end{example}

\section[typeTranslation]{Translation between classical types}

The word "type" is used to keep several parallel objects aligned:
root systems, Weyl groups, Lie groups, Hecke algebras, Schubert calculus, and
families of combinatorial objects.  The following table records the most
common translations.

\begin{array}{llll}
\toprule
\text{Type} & \text{Weyl group} & \text{Combinatorial model} & \text{Typical pages}\\
\midrule
A_{n-1}
  & \symS_n
  & \text{permutations}
  & \text{\hyperref[permutationMatrix]{permutations}, \hyperref[schubert]{Schubert}}\\
B_n
  & \symB_n
  & \text{signed permutations}
  & \text{\hyperref[permutationsTypeB]{type \(B\) permutations}}\\
C_n
  & \symB_n
  & \text{signed permutations}
  & \text{\hyperref[schurSymplectic]{symplectic Schur}}\\
D_n
  & \symD_n
  & \text{even signed permutations}
  & \text{\hyperref[weylLieGroups]{root systems}}\\
\widetilde A_{n-1}
  & \text{affine symmetric group}
  & \text{affine permutations}
  & \text{\hyperref[notationAffinePermutations]{affine permutations}}\\
\bottomrule
\end{array}

\svgimg[width=0.78\textwidth]{svg-images/coxeter-types-diagrams.svg}{
Small Coxeter diagrams for type \(A\), type \(B/C\), type \(D\), and affine
type \(A\).}

Types \(B_n\) and \(C_n\) have different crystallographic root systems but the
same Weyl group, the hyperoctahedral group.  This is why type \(B\) and type
\(C\) combinatorics often share signed-permutation models while differing in
root-length conventions, representation theory, or geometry.

In these diagrams an unlabeled edge denotes braid parameter \(m=3\), while
the doubled edge in type \(B/C\) records \(m=4\).  The diagram is often the
fastest translation layer between a presentation by generators and relations
and the corresponding root-system or Weyl-group model.

\section[coxeterCombinatorics]{Coxeter combinatorics}

Several constructions on this site have type \(A\) versions and then analogues
in other types:
\begin{itemize}
\item \hyperref[bruhatOrder]{Bruhat order} and weak order organize Schubert
calculus and reduced-word combinatorics.
\item \hyperref[heckeAlgebra]{Hecke algebras} deform group algebras of
Coxeter groups, while \hyperref[zeroHeckeAlgebra]{0-Hecke algebras} replace
the quadratic relation by an idempotent relation.
\item \hyperref[kazhdanLusztigPolynomials]{Kazhdan--Lusztig polynomials} are
indexed by intervals in Bruhat order.
\item \hyperref[cspWCatalan]{Coxeter--Catalan objects} generalize ordinary
Catalan objects from type \(A\) to all finite types.
\item \hyperref[rootSystem]{Root systems} provide the geometric language of
roots, coroots, chambers, and Weyl groups.
\end{itemize}

\section[coxeterTypeA]{Type \(A\)}

Type \(A_{n-1}\) is the symmetric group \(\symS_n\).  It is the ambient type
for ordinary Schur functions, Schubert polynomials, RSK, Young tableaux,
ordinary parking functions, and most of the basic examples on this site.

In type \(A\), the root system can be realized as
\[
  \{e_i-e_j : 1\leq i\ne j\leq n\},
\]
and the simple roots are \(e_i-e_{i+1}\).  The simple reflection \(s_i\)
swaps \(i\) and \(i+1\).

\section[coxeterTypeB]{Types \(B\) and \(C\)}

The Weyl group in types \(B_n\) and \(C_n\) is the group of
\hyperref[permutationsTypeB]{signed permutations}.  Combinatorially, this
means permutations of
\[
  \{\pm 1,\pm 2,\dotsc,\pm n\}
\]
with \(w(-i)=-w(i)\).  It has order \(2^n n!\).  Type \(B\) Eulerian
polynomials, type \(B\) noncrossing partitions, and type \(B\) parking-like
objects are typical examples of type \(A\) constructions with signed
analogues.

\section[coxeterTypeD]{Type \(D\)}

Type \(D_n\) is the subgroup of the signed permutation group consisting of
signed permutations with an even number of sign changes.  It has order
\(2^{n-1}n!\).  Type \(D\) often behaves like type \(B\) with an even-sign
condition, but this small-looking change can affect descents, Catalan objects,
and root-system conventions.

\section[coxeterAffine]{Affine types}

Affine Coxeter groups are infinite Coxeter groups obtained by adding an
affine simple reflection.  In type \(\widetilde A\), the resulting group can
be modeled by \hyperref[notationAffinePermutations]{affine permutations}.
Affine type \(A\) is the natural setting for affine Stanley symmetric
functions, \(k\)-Schur functions, affine Grassmannian combinatorics, and
related Hecke-algebra constructions.
